% ============================================================
% 论文写作素材（独立文件模板，供队友在阶段 4 写作取材）
% 用法：按本模板生成独立文件 solution/internal-reports/writing-material-<subN>.tex
%       （不 \input 进内部报告正文——正文与素材目的分离，见 TRAE-报告.md L1）
% 注意：本文件不进入终稿正文，仅作为写作交接材料；字段填写不全视为报告未定稿。
%       溯源/字段路径/可写禁写句一律放本文件，正文不承载。
% ============================================================
\documentclass{article}   % 独立编译用；若并入报告正文编译则删除本行与 \begin{document} 等
\usepackage[UTF8]{ctex}
\usepackage{xcolor}
\usepackage{booktabs}
\newcommand{\todo}[1]{\textcolor{red}{[TODO: #1]}}
\begin{document}
\section{论文写作素材（供队友写作）—— 本文件为独立素材，对应内部报告 iter-<序号>-sub<N>-<名>.tex}

\subsection{章节映射}

本报告对应终稿第 \textbf{第 N 章}「\textbf{章节标题}」。建议终稿结构如下：
\begin{itemize}
  \item §X.1 建模与求解 \quad $\leftarrow$ 本报告 §2 模型 + §3 数据要点
  \item §X.2 结果分析 \quad $\leftarrow$ 本报告 §4 结果 + §5 基线对比
\end{itemize}

\subsection{故事线（按段给出主题句与写作要点）}

\begin{enumerate}
  \item 【章首段】本问任务转述 + 输入规模 + 关键约束。要点：\todo{一句话转述赛题要求，给出任务数/区域数/时域等关键数字}。表述库参考：\todo{如 5.1-48（业务情境）+ 5.1-49（编号列表）}。内嵌摘录：\todo{若 parts 文件有更贴合的原文摘录，整段摘录于此并附出处，如 2021C1-摘要}
  \item 【建模段】为什么用该方法 + 模型结构（决策变量/目标/约束）。要点：\todo{说明方法选择理由与模型骨架，引用报告 §2 公式编号}。表述库参考：\todo{如 1.1-1（资源优化式）+ 2.2-20（双目标式）}。内嵌摘录：\todo{同上，附出处}
  \item 【结果段】核心数字 + 对比结论。要点：\todo{引用下方关键数字表，写明与基线/上一问的对比方向}。表述库参考：\todo{如 3.2-33（新旧方案增益对比）}。内嵌摘录：\todo{同上，附出处}
\end{enumerate}

\textbf{表述库内嵌规则}：队友无法直接访问 parts 文件（暂不共享），生成者必须把本问用到的 parts 原文摘录直接粘进「内嵌摘录」字段并注明出处（论文ID-章节），确保队友不依赖 parts 文件即可写作。

\subsection{关键数字表}

\begin{center}
\footnotesize
\begin{tabular}{llll}
\toprule
数字 & 含义 & 来源（脚本/文件） & 论文引用位置 \\
\midrule
\todo{2166.64M} & \todo{S4 主解运行成本} & \todo{sub4-model.py / iter-04 §4} & \todo{§7 结果段} \\
\bottomrule
\end{tabular}
\end{center}

\subsection{图表清单}

\begin{center}
\footnotesize
\begin{tabular}{lll}
\toprule
图/表 & 论文用途 & 建议插入位置 \\
\midrule
\todo{sub4-region-cost-carbon.pdf} & \todo{六区成本碳排分布} & \todo{§7 结果段} \\
\bottomrule
\end{tabular}
\end{center}

\subsection{口径与局限（可写句 / 禁写句）}

\begin{itemize}
  \item 可写：\todo{该子问题可以安全写进论文的结论句，如"成本降低 2.1\%，包含任务迁移全部收益"} 
  \item 禁止写：\todo{该子问题不得写进论文的过度承诺句，如"该收益为耦合增益"（口径声明见报告正文，不得删改）}
\end{itemize}

\subsection{AI 工具使用标注}

本子问题涉及 AI 使用记录：\todo{[AI-4-1] 等编号}，供阶段 3.3 AI 使用声明引用。

\end{document}
